\chapter{Conclusion}\label{chap:conclusion}

\section{Summary}
This thesis asked how retrieval quality affects hallucination in RAG with small open-weight language
models, and whether a rule that needs no training can reduce it. Two models answered
\val{data.eval_questions} multi-hop HotpotQA questions under six conditions, from no retrieval to perfect
retrieval, over a corpus of \val{data.passages} passages. A language-model judge labelled the answers as
supported, contradicted or unsupported, and it was checked against a second judge and the author's labels on 100
answers.

Retrieval quality varied widely: the share of all \val{data.questions} questions with both gold passages in the
top five ranged from \val{ret.bm25.both5.all} for BM25 to \val{ret.hybrid_rerank.both5.all} for hybrid retrieval with
reranking.

For RQ1, whether retrieval found every hop matters much more than which retriever was used. With both gold passages in
the top five, 7.0\,\% and 4.5\,\% of the answers of the two models are hallucinated; with only one, 63.1\,\% [resampling
questions: 48.4, 77.0] and 77.7\,\% [65.7, 88.3]. The gap holds when a question is compared with itself (64.1\,\% against
10.4\,\% and 78.3\,\% against 5.0\,\%) and when the author's labels replace the judge's (26 of 42 against 3 of 42). The
intervals of the four retrievers overlap. Agreement with the human check was only moderate ($\kappa=0.54$), and it suggests
that the absolute rates may be too high (about 19.1\,\% overall against 24.2\,\%).

For RQ2, retrieval confidence does not work: the top1 score is near chance (AUROC 0.46--0.54). A simple check of whether
the answer's words appear in the passages does better (AUROC 0.75--0.77 for \texttt{gpt-oss-20b}, 0.60--0.63 for
\texttt{qwen3.8-27b}). In its best case it lowers hallucination from 21.3\,\% to 11.0\,\% at 74.9\,\% coverage, at a cost in
correct answers, and it helps much less for the model that more often picks unsupported spans from the context.

\section{Limitations}\label{sec:limitations}
The conclusions are limited in several ways.
\begin{itemize}
\item \textbf{Few questions.} Generation and judging use \val{data.eval_questions} questions per condition,
  because of the free API limits, so intervals are wide and small differences may not be real. The abstention
  results rest on small test halves.
\item \textbf{One sample of questions.} Retrieval on the first \val{data.eval_questions} questions differs from
  retrieval on all \val{data.questions} (Section~\ref{sec:res-retrieval}), so the sample may not represent the full
  question set.
\item \textbf{An LLM as judge.} The hallucination labels come from a language model, which can make mistakes
  \cite{es2024ragas,niu2024ragtruth}, and every hallucination rate depends on them. The second judge shares a family with
  one generator. The human check covers 100 answers by one rater who knew the study's results, and agreement with it was only
  moderate. The causes were assigned by one rater who saw the judge's reason, so they are not independent.
\item \textbf{Faithfulness is not truth.} A supported answer is backed by the passages, which does not make it
  correct, and an unsupported answer may still be true.
\item \textbf{Association, not cause.} The link between gold passages found and hallucination is observed, not
  controlled, as questions that are hard to retrieve for may also be harder to answer. The same-question check reduces this
  concern, but it covers only questions where the retrievers disagree.
\item \textbf{Different context sizes.} The oracle condition shows the model two passages and the retrievers show
  five, so oracle results combine the effect of better evidence with that of a shorter context. Passage order also
  follows the retriever, and order can affect models \cite{liu2024lost}.
\item \textbf{One data set and two models.} HotpotQA is built from Wikipedia with short answers, and every
  question here is multi-hop. The results need not carry over to other domains, long answers or other models.
\item \textbf{Hosted models.} The models are reached through third-party APIs and may change over time.
\item \textbf{Abstention was applied after generation.} The rule is evaluated on saved answers, on test halves of about 49
  answers. The threshold needs labelled answers and is chosen for one objective, which weighs a correct answer and a
  hallucination equally. The lexical check suits short answers and may not carry over to long ones.
\item \textbf{Abstained answers are excluded} from the hallucination rate, and abstention rates differ between conditions (up
  to 11\,\%), which can slightly affect comparisons.
\end{itemize}

\section{Future Work}
Several directions follow from the results and the limitations.
\begin{itemize}
\item \textbf{More questions.} Running all \val{data.questions} questions would narrow the intervals and allow a
  larger test half for the abstention rule.
\item \textbf{Second-hop retrieval.} Since a missing hop goes with most hallucinations, retrieval that targets the second hop,
  for example iterative retrieval, should be tested directly.
\item \textbf{More confidence signals.} Signals that can catch unsupported spans from the context, which a lexical check cannot,
  could be compared with the signals tested here.
\item \textbf{Other prompts and context sizes.} Offering ``I don't know'' as an option in the prompt, forcing
  citations, or varying the number of passages would show how much can be gained without any retrieval signal.
\item \textbf{Broader settings.} Other data sets, longer answers and other models would test how general the
  findings are.
\end{itemize}
